\section{Complete results}
\label{appx:complete_results}



\subsection{TS-Memory across frozen TSFM backbones}
\label{appx:different-TSFMs-with-TS-Memory}

Table~\ref{tab:ts_plugmem_full} reports the full horizon-wise MSE/MAE results for TS-Memory on four frozen TSFM backbones (ChronosBolt, Chronos2, Sundial, and TimesFM), complementing the horizon-averaged summary in the main table. Across datasets and horizons, PlugMem yields consistent long-horizon gains without updating backbone parameters. In terms of MSE, average reductions are typically in the 5\%--7\% range for ChronosBolt, Chronos2, and Sundial, while TimesFM often sees larger reductions closer to 7\%--9\%. The complete per-dataset, per-horizon breakdown confirms that improvements are broad-based rather than driven by a small subset of benchmarks or horizons. The horizon-dependent patterns are not uniform across backbones, which motivates reporting the full table. Chronos2 shows a clearer tendency for gains to grow with the forecasting horizon, consistent with intuition that learned memory corrections can mitigate error accumulation during long iterative rollouts. Sundial and TimesFM, by contrast, often realize their strongest relative gains at short and medium horizons while remaining positive at $H=720$. This pattern suggests part of their long-horizon strength may already be captured by their native decoding behavior, with PlugMem providing complementary refinements.

The dataset-level view also reveals heterogeneity that can be obscured by horizon-averaged summaries. Weather and the minute-level ETT benchmarks repeatedly show larger relative improvements, whereas Traffic is more challenging for some backbones. Finally, MAE improvements are generally smaller than MSE improvements, indicating that TS-Memory primarily suppresses occasional large-error outliers rather than uniformly shifting all predictions.







\subsection{Comparison with adaptation baselines}
\label{appx:rag-TS-memory}

\noindent\textbf{Accuracy and calibration.}
Table~\ref{tab:chronosbolt_base_ablation} provides the full horizon-wise comparison on ChronosBolt against online retrieval baselines, including not only point errors (MSE/MAE) but also CRPS. TS-Memory achieves a stable MSE reduction of roughly 5\%--6\% across horizons, while its CRPS improvements become more pronounced as the horizon increases (from around 4\% at $H=96$ to about 8\% at $H=720$). This horizon-dependent CRPS trend is consistent with distillation being particularly beneficial for probabilistic quality when long-range rollouts amplify uncertainty. RAFT yields smaller and flatter gains (around a 3\% MSE reduction with about a 1\% MAE decrease), suggesting that retrieval helps but does not fully correct long-horizon drift. TS-RAG is less stable in the full breakdown: its average MSE reduction remains below about 2.2\% and MAE can degrade at some horizons, implying that directly injecting retrieved trajectories may introduce bias or noise even when it occasionally helps MSE. Dataset-level details reinforce the same message: the strongest distillation benefits concentrate on regimes with clearer repeating structure (e.g., Weather and minute-level ETT), while on Exchange the longest-horizon point-error gains can be small even when CRPS still improves, indicating better-calibrated predictive distributions even when the mean prediction saturates.


\noindent\textbf{Per-horizon latency decomposition.}
Table~\ref{tab:chronosbolt_time_latency} clarifies why online retrieval can be expensive in practice. For both RAFT and TS-RAG, retrieval introduces a substantial overhead and tends to increase with the prediction horizon, consistent with repeated retrieval during rollout or higher effective context cost at longer horizons. At $H=96$, retrieval already accounts for roughly one third of total latency for RAFT and more than half for TS-RAG, corresponding to total slowdowns of about 71\% and 193\% relative to the no-retrieval baseline. As the horizon grows to 720, the forward pass naturally becomes dominant so the retrieval fraction decreases; however, the overall overhead remains large at about 46\% for RAFT and about 126\% for TS-RAG. In contrast, TS-Memory removes retrieval latency entirely and introduces only a small forward-pass overhead (under 6\% across horizons), with the relative overhead becoming less visible for long horizons because decoding cost is unavoidable. Overall, among retrieval-based baselines, distillation is the only approach here that scales to long horizons without making retrieval a dominant runtime bottleneck.

\noindent\textbf{End-to-end evaluation time.}
When timing full evaluation runs rather than individual queries, Table~\ref{tab:chronosbolt_total_time} shows the same trend at wall-clock scale. Even if RAFT slightly reduces forward time in some settings, the added retrieval stage dominates the total runtime, yielding around 57\%--71\% slowdowns at short horizons and still about 40\% extra time at $H=720$. TS-RAG is substantially heavier, producing roughly 118\%--193\% slowdowns depending on horizon while keeping retrieval responsible for about half of the total runtime. TS-Memory remains close to the baseline with only 3.5\%--5.7\% extra total time, and this gap consistently shrinks at larger horizons because the forward computation grows. We note that occasional small speedups on some datasets are largely explained by fixed warm-up overhead and a smaller number of forecasting queries, which makes end-to-end timing less stable. These end-to-end results also highlight a dataset-size effect: on large-scale datasets such as Traffic, the number of forecasting calls is high, so retrieval overhead can accumulate into a prohibitive wall-clock penalty even when the per-query retrieval fraction appears moderate.




\subsection{Transfer generality}
\label{appx:transferability}

\noindent\textbf{Scaling across backbone sizes.}
Shrinking the backbone does not erase the benefit of retrieval distillation, as Table~\ref{tab:chronosbolt_ts_plugmem_full} demonstrates across ChronosBolt sizes. Averaged over horizons and datasets, TS-Memory typically delivers about 4.8\%--5.8\% MSE reduction across sizes, showing that gains are not limited to large backbones. The horizon-wise breakdown reveals an interaction between capacity and long-range forecasting: the mini model exhibits growing relative gains as the horizon increases (from about 4.5\% at $H=96$ to roughly 7.1\% at $H=720$), suggesting that smaller backbones benefit more from learned correction when compounding errors become severe. The tiny model also remains robust, maintaining improvements across all horizons. MAE follows similar but smaller trends, consistent with TS-Memory preferentially removing occasional large deviations rather than uniformly shifting all outputs.


\noindent\textbf{Cross-model transfer across backbones.}
Cross-model transfer in Table~\ref{tab:plugmem_merged} indicates that PlugMem generalizes beyond the specific teacher--target pairing used during distillation. Averaged over datasets and horizons, transfer yields around 7\% MSE reduction on ChronosBolt targets and around 6\% on Sundial targets, while TimesFM targets benefit the most with roughly 9\% MSE reduction. Gains persist across horizons, with TimesFM maintaining strong improvements from $H=96$ through 720, suggesting that the memory captures longer-range error dynamics rather than only short-horizon fixes. Teacher choice is also less critical than one might expect: Chronos2-PlugMem and Sundial-PlugMem produce comparable gains on shared targets, implying that TS-Memory distills a common retrieval signal rather than teacher-specific quirks. The per-dataset view further shows that transfer can be especially valuable under domain difficulty or shift (e.g., large gains on ETTm2 for TimesFM and on Exchange for Sundial), whereas Traffic often sees smaller changes concentrated in MSE rather than MAE.



\noindent\textbf{Transfer under domain shift.}
Table~\ref{tab:domain_split_lora_tsmemory} isolates how the \emph{train–test shift in retrieval supervision} impacts TS-Memory and LoRA: we keep the frozen backbone (ChronosBolt), retrieval hyperparameters, and training recipe fixed, varying only which split/domain provides supervision for the offline retrieval teacher (settings summarized in Table~\ref{tab:setting_test_relation}). A clear pattern emerges: domain alignment is the key driver, and the advantage grows with forecasting horizon. Across all configurations, average MSE reduction widens from $H{=}96$ to $H{=}720$, consistent with longer rollouts relying more heavily on retrieved analogs. Even cross-domain supervision provides measurable transfer, with gains increasing from about $1.4\%$ at $H{=}96$ to roughly $3.7\%$ at $H{=}720$, suggesting that generic temporal motifs can generalize beyond the source domain. Partially aligned domains (distribution-shift or multi-domain) roughly double the benefit, reaching about $4.5\%$--$4.8\%$ MSE reduction at $H{=}720$. In-domain supervision remains strongest, improving MSE by about $5.0\%$ at $H{=}96$ and nearly $7.8\%$ at $H{=}720$. In contrast, LoRA is substantially more sensitive to domain mismatch—it can underperform the frozen backbone under cross-domain training and shows less stable benefits when domains are misaligned. These results demonstrate that retrieval-distilled supervision provides a more robust transfer pathway than parameter-efficient fine-tuning: PlugMem preserves stronger gains under domain shifts by correcting phase/drift errors through external memory without re-fitting backbone parameters.






\noindent\textbf{Scaling PlugMem capacity.}
Following the analysis in Section~\ref{sec:model_analysis}, we evaluate three PlugMem variants (small, base, and large). Table~\ref{tab:mem-size} shows only minor differences across these configurations, with all three landing near a 5.5\% average MSE reduction and about 1.9\%--2.0\% MAE reduction across datasets and horizons. The horizon-level breakdown also indicates that bigger is not uniformly better: the base memory can be slightly stronger at shorter horizons, the large memory can edge out at some medium horizons, and the ordering can flip across horizons, suggesting a non-monotonic interaction between memory capacity and dataset dynamics rather than a simple scaling law. Dataset-level behavior points to the same conclusion: memory size has limited impact on Traffic where improvements remain small, whereas on Weather and minute-level ETT data even the smallest memory delivers substantial gains. Overall, most of the retrieval-distillation benefit is captured with only a few million parameters, making PlugMem tunable for parameter efficiency without sacrificing much accuracy.









